% Part: computability
% Chapter: computability-theory
% Section: fixed-point-thm

\documentclass[../../../include/open-logic-section]{subfiles}

\begin{document}

\olfileid[pt]{cmp}{thy}{fix}
\olsection{O teorema do ponto fixo}

Consideremos novamente o problema da parada. Como notação temporária,
escrevamos $\gn{\cfind{x}(y)}$ para $\tuple{x, y}$; pense nisso como
uma representação de um ``nome'' para o valor $\cfind{x}(y)$. Com
essa notação, podemos reformular uma de nossas demonstrações de que o
problema da parada é indecidível.

Pergunta: existe uma função computável $h$ com a
seguinte propriedade? Para todos $x$ e $y$,
\[
h(\gn{\cfind{x}(y)}) =
\begin{cases}
1 & \text{se $\cfind{x}(y) \fdefined$} \\
0 & \text{caso contrário.}
\end{cases}
\]

Resposta: Não; caso contrário, a função parcial
\[
g(x) \simeq
\begin{cases}
0 & \text{se $h(\gn{\cfind{x}(x)}) = 0$} \\
\text{indefinida} & \text{caso contrário}
\end{cases}
\]
seria computável e, portanto, teria algum índice~$e$. Mas então teríamos
\[
\cfind{e}(e) \simeq
\begin{cases}
0 & \text{se $h(\gn{\cfind{e}(e)}) = 0$} \\
\text{indefinida} & \text{caso contrário,}
\end{cases}
\]
e, nesse caso, $\cfind{e}(e)$ estaria definida se, e somente se, não
estivesse, uma contradição.

Agora examine a equação com $\cfind{e}$. Em certo sentido, há nela um
caso de autorreferência: fizemos com que o valor de $\cfind{e}(e)$
dependesse de $\gn{\cfind{e}(e)}$ de determinada maneira. O teorema do
ponto fixo diz que, em geral, {\em podemos} fazer isso---e não apenas
para demonstrar contradições.

A \olref{lem:fixed-equiv} apresenta duas maneiras equivalentes de
enunciar o teorema do ponto fixo. Do ponto de vista lógico, a
equivalência dos enunciados decorre do fato de ambos serem verdadeiros;
mas o que realmente queremos dizer é que cada um decorre diretamente
do outro, de modo que podem ser tomados como enunciados alternativos
do mesmo teorema.

\begin{lem}
\ollabel{lem:fixed-equiv}
Os seguintes enunciados são equivalentes:
\begin{enumerate}
\item Para toda função parcialmente computável $g(x,y)$, existe um
  índice~$e$ tal que, para todo~$y$,
\[
\cfind{e}(y) \simeq g(e,y).
\]
\item Para toda função computável~$f(x)$, existe um índice~$e$ tal
  que, para todo~$y$,
\[
\cfind{e}(y) \simeq \cfind{f(e)}(y).
\]
\end{enumerate}
\end{lem}

\begin{proof}
$(1) \Rightarrow (2)$: Dada $f$, defina $g$ por
$g(x,y) \simeq \fn{Un}(f(x),y)$. Use (1) para obter um índice~$e$ tal
que, para todo~$y$,
\begin{align*}
\cfind{e}(y) & = \fn{Un}(f(e),y) \\
& = \cfind{f(e)}(y).
\end{align*}

$(2) \Rightarrow (1)$: Dada $g$, use o teorema $s$-$m$-$n$ para obter
$f$ tal que, para todos $x$ e~$y$,
$\cfind{f(x)}(y) \simeq g(x,y)$. Use (2) para obter um índice~$e$ tal
que
\begin{align*}
\cfind{e}(y) & = \cfind{f(e)}(y) \\
& = g(e,y).
\end{align*}
Isso conclui a demonstração.
\end{proof}

\begin{explain}
Antes de mostrar que o enunciado (1) é verdadeiro (e, portanto, (2)
também), considere como ele é estranho. Pense em $e$ como um programa
de computador; o enunciado (1) diz que, dada qualquer função
parcialmente computável $g(x,y)$, você pode encontrar um programa de
computador $e$ que calcula $g_e(y) \simeq g(e,y)$. Em outras palavras,
você pode encontrar um programa de computador que calcula uma função
que faz referência ao próprio programa.
\end{explain}

\begin{thm}
Os dois enunciados da \olref{lem:fixed-equiv} são
verdadeiros. Especificamente, para toda função parcialmente computável
$g(x,y)$, existe um índice~$e$ tal que, para todo~$y$,
\[
\cfind{e}(y) \simeq g(e,y).
\]
\end{thm}

\begin{proof}
Os ingredientes já estão implícitos na discussão do problema da parada
acima. Seja $\fn{diag}(x)$ uma função computável que, para cada $x$,
devolve um índice da função
$f_x(y) \simeq \cfind{x}(x,y)$, isto é,
\[
\cfind{\fn{diag}(x)}(y) \simeq \cfind{x}(x,y).
\]
Pense em $\fn{diag}$ como uma função que transforma um programa para
uma função binária em um programa para uma função unária, obtida
fixando o programa original como seu primeiro argumento. A função
$\fn{diag}$ pode ser definida formalmente da seguinte maneira:
primeiro, defina $s$ por
\[
s(x,y) \simeq \fn{Un}^2(x,x,y),
\]
em que $\fn{Un}^2$ é uma função ternária universal para funções binárias
parcialmente computáveis. Então, pelo teorema $s$-$m$-$n$, podemos
encontrar uma função recursiva primitiva $\fn{diag}$ que satisfaça
\[
\cfind{\fn{diag}(x)}(y) \simeq s(x,y).
\]

Agora, defina a função $l$ por
\[
l(x,y) \simeq g(\fn{diag}(x),y).
\]
e seja $\gn{l}$ um índice de $l$. Por fim, seja
$e = \fn{diag}(\gn{l})$. Então, para todo $y$, temos
\begin{align*}
\cfind{e}(y) & \simeq \cfind{\fn{diag}(\gn{l})}(y) \\
& \simeq \cfind{\gn{l}}(\gn{l}, y) \\
& \simeq l(\gn{l}, y) \\
& \simeq g(\fn{diag}(\gn{l}),y) \\
& \simeq g(e, y),
\end{align*}
como queríamos.
\end{proof}

\begin{explain}
O que está acontecendo? Suponha que você receba a tarefa de escrever
um programa de computador que imprima a si próprio. Suponha, ainda,
que esteja trabalhando com uma linguagem de programação dotada de uma
biblioteca rica e estranha de funções sobre cadeias. Em particular,
suponha que sua linguagem de programação tenha uma função $\fn{diag}$
que opere da seguinte maneira: dada uma cadeia de entrada~$s$,
$\fn{diag}$ localiza cada ocorrência do símbolo `x' em~$s$ e a substitui
por uma versão entre aspas da cadeia original. Por exemplo, dada a
cadeia
\begin{quote}
\begin{verbatim}
olá x mundo
\end{verbatim}
\end{quote}
como entrada, a função devolve
\begin{quote}
\begin{verbatim}
olá 'olá x mundo' mundo
\end{verbatim}
\end{quote}
como saída. Nesse caso, é fácil escrever o programa desejado; você
pode verificar que
\begin{quote}
\begin{verbatim}
print(diag('print(diag(x))'))
\end{verbatim}
\end{quote}
resolve o problema. Para linguagens de programação mais comuns, como
C++ e Java, a mesma ideia (com uma implementação mais complexa) ainda
funciona.

Estamos a apenas alguns passos da demonstração do teorema do ponto
fixo. Suponha que uma variante da função de impressão
$\fn{print}(x,y)$ aceite uma cadeia $x$ e outro argumento numérico $y$,
e imprima repetidamente a cadeia $x$, $y$ vezes. Então o ``programa''
\begin{quote}
\begin{verbatim}
getinput(y);
print(diag('getinput(y); print(diag(x), y)'), y)
\end{verbatim}
\end{quote}
imprime a si próprio $y$ vezes na entrada $y$. Substituir o esqueleto
$\fn{getinput}$---$\fn{print}$---$\fn{diag}$ por uma função arbitrária
$g(x,y)$ produz
\begin{quote}
\begin{verbatim}
g(diag('g(diag(x), y)'), y)
\end{verbatim}
\end{quote}
que é um programa que, na entrada $y$, executa $g$ sobre o próprio
programa e sobre $y$. Entendendo ``colocar entre aspas'' como ``usar
um índice para'', obtemos a demonstração acima.

Por enquanto, não há problema.\ Se você quiser considerar a demonstração
um artifício formal ou magia negra. Mas você deve ser capaz de
reconstruir os detalhes do argumento apresentado acima. Quando
demonstrarmos os teoremas da incompletude (e o ``teorema do ponto
fixo'' relacionado), discutiremos outras maneiras de entender por que
ele funciona.
\end{explain}

\begin{tagblock}{lambda}
\begin{digress}
A mesma ideia pode ser usada para obter um combinador de ``ponto fixo''.
Suponha que você tenha um termo lambda $g$ e queira outro termo $k$ com
a propriedade de que $k$ é $\beta$-equivalente a $gk$. Defina os termos
\[
\fn{diag}(x) = xx
\]
e
\[
l(x) = g(\fn{diag}(x))
\]
usando nossas convenções de notação; em outras palavras, $l$ é o termo
$\lambd[x][g(xx)]$. Seja $k$ o termo $ll$. Então temos
\begin{align*}
k & = (\lambd[x][g(xx)])(\lambd[x][g(xx)]) \\
& \red  g((\lambd[x][g(xx)])(\lambd[x][g(xx)])) \\
& = gk.
\end{align*}
Se tomarmos
\[
Y = \lambd[g][((\lambd[x][g(xx)])(\lambd[x][g(xx)]))]
\]
então $Yg$ e $g(Yg)$ se reduzem a um termo comum; logo,
$Yg \equiv_\beta g(Yg)$. Isso é conhecido como ``combinador de
Curry''. Se, em vez disso, tomarmos
\[
Y = (\lambd[xg][g(xxg)])(\lambd[xg][g(xxg)])
\]
então, de fato, $Yg$ se reduz a $g(Yg)$, o que é uma afirmação mais
forte. Essa última versão de $Y$ é conhecida como ``combinador de
Turing''.
\end{digress}
\end{tagblock}

\end{document}
